Technológie umelých neurónových sietí a~hlbokého učenia zohrávajú čoraz dôležitejšiu úlohu v~oblasti vnímania a~lokalizácie autonómnych vozidiel. Oblasť fúzie senzorických dát nepatrí medzi tradičné príklady implementácie metód umelej inteligencie, vývoj v~nej sa dlho uberal skôr analytickými a~pravdepodobnostnými smermi. Problém s~náročnosťou modelovania a~analytického či pravdepodobnostného spracovania senzorických dát v~súlade s~rozvojom umelej inteligencie podnecuje rozvoj aplikácií aj v~týchto kombináciách\@. \acrshort{uns} majú potenciál zohrať úlohu nie len v~samotnom spájaní senzorických tokov do jedného, ale najmä v~inteligentnom predspracovaní senzorických dát za účelom potlačenia šumov či iných nežiadúcich faktorov. Zlepšenie lokalizácie by bolo možné tiež dosiahnuť detekciou výrazných orientačných bodov (napr.\ značky, stĺpy v~exteriéri alebo dvere, schodiská v~interiéri) a~ich porovnávaním s~mapou, pričom využitie 3D neurónových sietí zvyšuje robustnosť voči rušivým vplyvom prostredia. Kým konvolučné siete sú dnes štandardom pri spracovaní obrazov, v~lokalizácii a~mapovaní je potenciál iných architektúr, najmä rekurentných, stále nevyužitý. Budúci výskum sa preto môže zamerať na vytváranie komplexnejších riešení a~zároveň posilniť spoľahlivosť, opakovateľnosť a~odolnosť voči útokom kombinovaním dátovo orientovaných prístupov s~klasickými modelmi.